\documentclass[10pt,a4paper]{amsart}
\usepackage[margin=19mm]{geometry}
\usepackage{amsmath,amssymb,mathtools}
\usepackage[scr=rsfs,cal=euler]{mathalfa}
\usepackage{xfrac}
\usepackage[unicode,hidelinks,bookmarks=false]{hyperref}
\newcommand{\ie}{i.e.\ }
\newcommand{\cf}{cf.\ }
\newcommand{\resp}{respectively\ }
\newcommand{\Subsection}{\subsection}
\providecommand{\nobreakdash}{\nobreak\hbox{-}}
\setlength{\emergencystretch}{2em}
\multlinegap0pt
\usepackage{amssymb}
\usepackage[shortlabels]{enumitem}
\usepackage{xcolor}
\usepackage{tikz}
\usepackage{float}
\newtheorem{theorem}{Theorem}
\newtheorem{lemma}{Lemma}
\title{On the Cauchy-Hamel Continuity of Real Functions}
\author{Gabriel Istrate}
\date{Source: arXiv:2608.16942v2 (CC BY 4.0)}
\begin{document}
\maketitle

\noindent\emph{Source and licence.} On the Cauchy-Hamel Continuity of Real Functions. Version arXiv:2608.16942v2 (CC BY 4.0), distributed by arXiv under the Creative Commons Attribution 4.0 International licence, http://creativecommons.org/licenses/by/4.0/, as recorded in the arXiv licence metadata for this exact version and shown on the abstract page https://arxiv.org/abs/2608.16942v2. Full licence notice: \texttt{licenses/arxiv-2608.16942-CC-BY-4.0.txt}.\par\smallskip

\noindent\textit{Editorial scope.} Counted unit: the article's theorem thm:cons2 (source0.tex lines 1886--1896). The article's complete original proof of it is delivered with it (source0.tex lines 1898--2036, one \texttt{proof} environment of 139 lines). No mathematics is written, repaired or re-derived: the statement and the proof are the article's own lines, in the article's own notation, and no retained line is substituted or re-flowed.  No further source-local context is needed: every cross-reference of every retained line resolves inside the two delivered spans.  Nothing was omitted from any retained span, so the package carries no omission record.  No retained line contains a citation, so the wrapper carries no bibliography and no bibliography entry.  No retained line contains a figure environment, an image-inclusion command or drawing code, so no image is delivered and the package declares no asset dependency.  Editorial formatting only, in addition: an AMS \texttt{amsart} preamble that restates the article's own declarations and macros used by the retained lines, namely the environments \texttt{theorem}, \texttt{lemma} on separate counters, together with the standard packages those lines need.  The retained environments print in this document as Theorem 1, Lemma 1 in order of appearance, whereas the article prints the same items with the numbering of its own full document.  The complete native source of the article as distributed in the arXiv e-print for this exact version is bundled unchanged as \texttt{sources/207864/source0.tex}.
\begin{theorem}
Let $A\subseteq\mathbb{Q}$ be finite, and let $f:A\rightarrow\mathbb{R}$ be arbitrary. Then one can
construct an extension $g:\mathbb{Q}\rightarrow\mathbb{R}$ such that:
\begin{enumerate}
\item[(1)] $g$ is continuous on $\mathbb{Q}$ (with the subspace topology inherited from $\mathbb{R}$),
\item[(2)] $g|_A = f$,
\item[(3)] there is no function $G:\mathbb{R}\rightarrow\mathbb{R}$ with $G|_{\mathbb{Q}}=g$ that is continuous on some
nondegenerate interval.
\end{enumerate}
\label{thm:cons2} 
\end{theorem}

\begin{proof}
Fix a countable dense set of irrationals. One explicit choice is as
follows: enumerate $\mathbb{Q}$ as $(r_k)_{k\ge 1}$, and set $\alpha_{k,m} = r_k + \frac{\sqrt{2}}{m} \quad (k,m\in\mathbb{N}, m\ge 1).$
Then each $\alpha_{k,m}\notin\mathbb{Q}$, and the set $D = \{\alpha_{k,m} : k,m\in\mathbb{N}, m\ge 1\} \subseteq \mathbb{R}\setminus\mathbb{Q}$
is dense in $\mathbb{R}$, because for any nonempty open interval $I$ one can pick a rational
$r_k\in I$ and then choose $m$ large enough so that $r_k+\sqrt{2}/m\in I$.

The following is a key observation: To prevent the existence of an interval on which some extension is continuous, it suffices to ensure that \empty{every} nondegenerate interval contains a point $x$
at which the rational limit of $g$ fails to exist. 

\begin{lemma} 
Let $g:\mathbb{Q}\rightarrow\mathbb{R}$ and let $F:\mathbb{R}\rightarrow\mathbb{R}$ satisfy $F|_{\mathbb{Q}}=g$. If $F$ is
continuous at some $x\in\mathbb{R}$, then
$ \lim_{\substack{q\rightarrow x\\ q\in\mathbb{Q}}} g(q)$
exists in $\mathbb{R}$ and equals $F(x)$. In particular, if the above rational limit does not exist
at $x$, then no extension $F$ can be continuous at $x$.
\label{lemma:extension}
\end{lemma}

\begin{proof}
Assume $F$ is continuous at $x$. Let $(q_n)\subseteq\mathbb{Q}$ with $q_n\rightarrow x$. Then $g(q_n)=F(q_n)\rightarrow F(x)$ by continuity of $F$ at $x$. Hence every rational sequence converging
to $x$ has image converging to the same limit $F(x)$, i.e. the rational limit exists
and equals $F(x)$.
\end{proof}



Now let $\emptyset \neq A = \{a_1 < \dots < a_n\} \subseteq \mathbb{Q}$ be finite and let $f:A\rightarrow\mathbb{R}$ be arbitrary. We construct the desired function $g$ in several steps.  
%\begin{itemize} 
%\item[-] \noindent \textbf{Step 1: construct a continuous real function $g_0$ matching $f$ on $A$.} 
First, define $g_0:\mathbb{R}\rightarrow\mathbb{R}$ by
piecewise linear interpolation:
\begin{itemize}
\item $g_0(x) = f(a_1)$ for $x\le a_1$;
\item for each $i=1,\dots,n-1$ and $x\in[a_i, a_{i+1}]$,
\begin{equation} g_0(x) = f(a_i) + \frac{x-a_i}{a_{i+1}-a_i} (f(a_{i+1})-f(a_i)); 
\end{equation} 
\item $g_0(x) = f(a_n)$ for $x\ge a_n$.
\end{itemize}
By construction $g_0$ is continuous on $\mathbb{R}$ and satisfies $g_0(a_i) = f(a_i)$ for all $i$.


%\item[-] \noindent \textbf{Step 2: construct a continuous ``oscillation'' term on $\mathbb{Q}$.} 

Next, enumerate the dense irrational set $D$ as $(\beta_{j})_{j\ge1}$ (any enumeration will do). For $j\ge1$ define
\begin{equation} 
\varphi_{j}(x)=\sin\left(\frac{1}{x-\beta_{j}}\right), \quad x\in\mathbb{R}\setminus\{\beta_{j}\}. \end{equation} 

Since $\beta_{j}\notin\mathbb{Q}$ the restriction $\varphi_{j}|_{\mathbb{Q}}$ is a well-defined continuous function on $\mathbb{Q}$ (it is the restriction of a continuous function on $\mathbb{R}\setminus\{\beta_{j}\}$, and no point of $\mathbb{Q}$ is removed). Define, for $q\in\mathbb{Q}$,
\begin{equation} 
H(q)=\sum_{j=1}^{\infty}2^{-j}\varphi_{j}(q). 
\end{equation} 
Since $|\varphi_{j}(q)|\le1$ for all $q\in\mathbb{Q}$, the Weierstrass M-test implies that the series converges uniformly on $\mathbb{Q}$, hence $H:\mathbb{Q}\rightarrow\mathbb{R}$ is continuous.

\vspace{0.2cm}
%\item[-] \noindent \textbf{Step 3: define a polynomial that is identically zero on $A$.} 

Now, let $P(x)=\prod_{a\in A}(x-a).$ $P$ is a polynomial with rational coefficients, so $P|_{\mathbb{Q}}$ is continuous on $\mathbb{Q}$ and $P(a)=0$ for all $a\in A$. 
%\item[-]\noindent \textbf{Step 4: Definition of $g$:} 

Finally define $g:\mathbb{Q}\rightarrow\mathbb{R}$ by
\begin{equation} g(q)=g_{0}(q)+P(q)H(q), \quad q\in\mathbb{Q}. 
\label{9.4} 
\end{equation} 
We can complete equation~(\ref{9.4}) to the case $A=\emptyset$ by taking in this case
\begin{equation} 
g_0\equiv 0\mbox{ and }P\equiv 1. 
\end{equation} 
%\end{itemize} 

\begin{lemma}
\label{thm:5.5}
The function $g:\mathbb{Q}\rightarrow\mathbb{R}$ constructed above satisfies:
\begin{enumerate}
    \item[(1)] $g$ is continuous on $\mathbb{Q}$,
    \item[(2)] $g|_{A}=f$,
    \item[(3)] no extension $F:\mathbb{R}\rightarrow\mathbb{R}$ with $F|_{\mathbb{Q}}=g$ is continuous on any nondegenerate interval.
\end{enumerate}
\end{lemma}


\begin{proof}
(1) The restriction $g_{0}|_{\mathbb{Q}}$ is continuous on $\mathbb{Q}$ because $g_{0}$ is continuous on $\mathbb{R}$. Also, $P|_{\mathbb{Q}}$ is continuous on $\mathbb{Q}$ and $H$ is continuous on $\mathbb{Q}$ by uniform convergence. Therefore the function $q\mapsto P(q)H(q)$ is continuous on $\mathbb{Q}$, and so is $g$.

(2) If $a\in A$ then $P(a)=0$, hence $g(a)=f(a).$

(3) Let $I\subseteq\mathbb{R}$ be a nondegenerate open interval. Since $D$ is dense, pick $j$ such that $\beta_{j}\in I$. We claim that the rational limit $\lim_{q\rightarrow\beta_{j},q\in\mathbb{Q}}g(q)$ does not exist. Write $H(q)=2^{-j}\varphi_{j}(q)+\sum_{m\ne j}2^{-m}\varphi_{m}(q).$
Define
\begin{equation} 
R(q):=\sum_{m\ne j}2^{-m}\varphi_{m}(q) \quad (q\in\mathbb{Q}), \end{equation}
and note that for each $m\ne j$ the value $\varphi_{m}(\beta_{j})$ is well-defined because $\beta_{j}\ne\beta_{m}$. Set $L:=\sum_{m\ne j}2^{-m}\varphi_{m}(\beta_{j}),$ 
which converges absolutely since $|\varphi_{m}(\beta_{j})|\le1$.
\begin{lemma}  We have:
\begin{equation} 
\lim_{q\rightarrow\beta_{j}}R(q)=L.
\end{equation} 
\end{lemma} 
\begin{proof} 
Indeed, fix $\epsilon>0$ and choose $N\geq j$ so that $\sum_{m>N}2^{-m}<\epsilon/4$. By continuity of each $\varphi_{m}$ at $\beta_{j}$ for $m\in\{1,...,N\}\setminus\{j\}$, there exists $\delta>0$ such that whenever $q\in\mathbb{Q}$ and $|q-\beta_{j}|<\delta$:
\begin{equation} 
\sum_{\stackrel{1\le m\le N}{m\neq j}}2^{-m}|\varphi_{m}(q)-\varphi_{m}(\beta_{j})|<\epsilon/2. 
\end{equation} 
For the tail we use $|\varphi_{m}|\le1$ to get
\begin{equation} 
\sum_{m>N}2^{-m}|\varphi_{m}(q)-\varphi_{m}(\beta_{j})|\le2\sum_{m>N}2^{-m}<\epsilon/2.
\end{equation}  

Combining these inequalities yields $|R(q)-L|<\epsilon$ for all $q\in\mathbb{Q}$ with $|q-\beta_{j}|<\delta$, proving the lemma.
\end{proof} 

Crucially, $P(\beta_{j})\ne0$ because $P$ has only rational roots (namely the elements of $A$), whereas $\beta_{j}$ is irrational.

Now choose two rational sequences $(q_{n})$ and $(q_{n}^{\prime})$ with $q_{n}\rightarrow\beta_{j}$ and $q_{n}^{\prime}\rightarrow\beta_{j}$ such that
\begin{equation} 
\varphi_{j}(q_{n})=\sin\left(\frac{1}{q_{n}-\beta_{j}}\right)\rightarrow1, \mbox{ and }
\end{equation} 
\begin{equation} 
\varphi_{j}(q_{n}^{\prime})=\sin\left(\frac{1}{q_{n}^{\prime}-\beta_{j}}\right)\rightarrow-1. 
\end{equation} 
Such sequences exist because $\sin(1/t)$ oscillates between $-1$ and $1$ arbitrarily close to $t=0$, and rationals are dense.

We have 
\begin{equation} 
g(q_{n})=g_{0}(q_{n})+P(q_{n})\left(2^{-j}\varphi_{j}(q_{n})+\sum_{m\ne j}2^{-m}\varphi_{m}(q_{n})\right)
\end{equation} 
hence 
\begin{equation} 
\lim_{n\rightarrow \infty} g(q_n)= g_{0}(\beta_{j})+P(\beta_{j})(2^{-j}\cdot1+L)
\end{equation} 
Similarly, it holds that 
\begin{equation} 
\lim_{n\rightarrow \infty} g(q_{n}^{\prime})= g_{0}(\beta_{j})+P(\beta_{j})(2^{-j}\cdot(-1)+L). 
\end{equation} 
These two limits differ by $2\cdot2^{-j}P(\beta_{j})\ne0$. Hence the rational limit of $g$ at $\beta_{j}$ does not exist.
 

By Lemma~\ref{lemma:extension}, no extension $F$ of $g$ to $\mathbb{R}$ can be continuous at $\beta_{j}$. Since $\beta_{j}\in I$, it follows that no such $F$ can be continuous on the interval $I$.
\end{proof}
\end{proof} 

\end{document}
